\documentclass[12pt,a4paper]{article}

\usepackage[T1]{fontenc}
\usepackage[utf8]{inputenc}
\usepackage[ngerman]{babel}
\usepackage{lmodern}
\usepackage{microtype}
\usepackage{geometry}
\geometry{a4paper, margin=2.5cm}

\usepackage{setspace}

\usepackage{graphicx}
\usepackage{booktabs}
\usepackage{longtable}
\usepackage{array}
\usepackage{tabularx}
\usepackage{multirow}
\usepackage{xcolor}
\usepackage{enumitem}
\usepackage{amsmath}
\usepackage{amssymb}

\usepackage[hidelinks]{hyperref}
\usepackage{cleveref}

\usepackage{titlesec}

\usepackage{abstract}
\renewcommand{\abstractnamefont}{\normalfont\bfseries\large}

\usepackage{tcolorbox}
\tcbuselibrary{skins,breakable}

\newtcolorbox{infobox}[1][]{
  colback=gray!10,
  colframe=gray!50,
  fonttitle=\bfseries,
  breakable,
  #1
}

% -------------------------------------------------------
% Metadaten
% -------------------------------------------------------
\title{
  {\bfseries Neuronale Netze und ISO 26262}\\
  {\Large Eine kritische Analyse der Vereinbarkeit unter besonderer Ber\"ucksichtigung von Zuverl\"assigkeit, Vertrauensw\"urdigkeit, Nachvollziehbarkeit und deterministischem Verhalten}}

\author{Stephan Epp}
\date{30. März 2026}

% -------------------------------------------------------
% Dokument
% -------------------------------------------------------
\begin{document}

\maketitle
\tableofcontents
\newpage

% ============================================================
\section{Einleitung}
\label{sec:einleitung}
% ============================================================
Der Einsatz von k\"unstlichen neuronalen Netzen (KNN) in sicherheitskritischen
elektronischen Steuerger\"aten (ECUs) des Automobils erfreut sich wachsender
Beliebtheit, insbesondere in der Realisierung virtueller Sensoren sowie bei
Fahrerassistenzfunktionen. Diese Entwicklung steht jedoch in einem grunds\"atzlichen
Spannungsverh\"altnis zum international g\"ultigen Sicherheitsentwicklungsstandard
ISO~26262, der deterministische, formal verifizierbare und l\"uckenlos r\"uckverfolgbare
Systeme fordert. Die vorliegende Arbeit analysiert diesen Zielkonflikt systematisch und
kommt zu dem Schluss, dass neuronale Netze in ihrer gegenw\"artigen Form die zentralen
Anforderungen der ISO~26262 -- insbesondere Determinismus, Nachvollziehbarkeit,
Verifikationsf\"ahigkeit sowie Anforderungsverfolgbarkeit -- strukturell nicht erf\"ullen
k\"onnen. Es wird argumentiert, dass systemarchitektonische
Kompensationsma\ss{}nahmen diese fundamental-epistemischen Defizite zwar
einschr\"anken, jedoch nicht beseitigen k\"onnen. Auf Basis einer eingehenden
Analyse des Normtextes, der Sicherheitsmethodik und des aktuellen Stands der
KI-Forschung wird ein differenziertes Fazit formuliert, das die strikte
Beschr\"ankung des KNN-Einsatzes auf unkritische Qualit\"atsmanagement-Funktionen
empfiehlt und den weiteren Einsatz ablehnt, solange keine formal gleichwertigen
Nachweisverfahren existieren.

Die Automobilindustrie befindet sich in einem tiefgreifenden technologischen Wandel.
Hochautomatisiertes Fahren, Elektromobilit\"at und die zunehmende Vernetzung von
Fahrzeugen mit ihrer Umgebung erfordern immer komplexere elektronische Steuersysteme.
In diesem Kontext gewinnen Methoden der k\"unstlichen Intelligenz (KI) und insbesondere
k\"unstliche neuronale Netze (KNN) an Bedeutung, da sie in der Lage sind, komplexe,
nichtlineare Zusammenh\"ange aus Daten zu erlernen -- F\"ahigkeiten, die klassischen,
regelbasierten Ans\"atzen oft verschlossen bleiben \cite{LeCun2015}.

Eine prominente Anwendung ist der sogenannte \emph{virtuelle Sensor}: Anstelle eines
kostenintensiven oder unter Serienbedingungen unpraktikablen physikalischen Sensors
soll ein KNN ausgehend von anderen, bereits vorhandenen Messgr\"o\ss{}en die gesuchte
physikalische Gr\"o\ss{}e sch\"atzen. Derartige Konzepte werden heute aktiv von
Werkzeuganbietern wie dSPACE mit ihrer TargetLink-Umgebung beworben und versprechen,
den Entwicklungsaufwand zu reduzieren und gleichzeitig die Leistungsf\"ahigkeit von
Steuerger\"aten zu erh\"ohen.

Diese Entwicklung wirft jedoch eine fundamentale Frage auf, die in der Industrie und
Wissenschaft bislang nicht befriedigend beantwortet wurde: \emph{K\"onnen neuronale
Netze die Anforderungen erf\"ullen, die der internationale Sicherheitsstandard
ISO~26262 an Software in sicherheitskritischen Steuerger\"aten stellt?}

ISO~26262, der Ma\ss{}gebliche Standard f\"ur die funktionale Sicherheit von
Stra\ss{}enfahrzeugen, definiert einen umfassenden Entwicklungsprozess, der auf
Determinismus, vollst\"andiger Anforderungsverfolgbarkeit, formaler Verifikation und
einer nachvollziehbaren, \"uberpr\"ufbaren Systemstruktur beruht \cite{iso26262_1}.
Diese Eigenschaften sind keine willk\"urlichen b\"urokratischen H\"urden, sondern
Konsequenz jahrzehntelanger Erfahrung mit sicherheitskritischer Systementwicklung und
dem Verst\"andnis, dass menschliches Leben von der zuverl\"assigen und vorhersagbaren
Funktion dieser Systeme abh\"angt \cite{Leveson2011}.

Die vorliegende Arbeit nimmt eine deutlich kritische Position gegen\"uber dem Einsatz
von neuronalen Netzen in sicherheitskritischen ECUs ein. Sie argumentiert, dass die
grunds\"atzliche Arbeitsweise neuronaler Netze -- das statistische Erlernen von
Mustern aus Daten ohne explizit programmierte Logik -- strukturell unvereinbar mit den
Kernanforderungen von ISO~26262 ist. Der Beitrag gliedert sich wie folgt:
\Cref{sec:grundlagen} legt die technischen und normativen Grundlagen. \Cref{sec:iso}
analysiert die relevanten ISO-26262-Anforderungen im Detail. \Cref{sec:knn_probleme}
identifiziert die spezifischen Defizite neuronaler Netze. \Cref{sec:argumente}
formuliert die zentralen Argumente gegen den Einsatz. \Cref{sec:gegenargumente}
setzt sich mit m\"oglichen Gegenargumenten und Kompensationsans\"atzen auseinander und
zeigt deren Grenzen auf. \Cref{sec:alternativen} diskutiert sichere Alternativen, und
\Cref{sec:fazit} schlie\ss{}t mit einem Fazit und Empfehlungen.

\textbf{Schl\"usselw\"orter:} ISO~26262, Funktionale Sicherheit, Neuronale
Netze, K\"unstliche Intelligenz, ASIL, Determinismus, Verifikation,
Anforderungsverfolgbarkeit, Virtuelle Sensoren, ECU

% ============================================================
\section{Grundlagen}
\label{sec:grundlagen}
% ============================================================

\subsection{Funktionale Sicherheit und ISO~26262}
\label{subsec:iso_grundlagen}

Funktionale Sicherheit bezeichnet die Eigenschaft eines Systems, auf fehlerhafte
Zust\"ande in einer f\"ur die Umgebung sicheren Weise zu reagieren
\cite{iec61508}. Im automobilen Kontext ist ISO~26262 die ma\ss{}gebliche Norm. Sie
entstand als Ableitung der allgemeinen Norm IEC~61508 und wurde erstmals 2011
ver\"offentlicht; die zweite, derzeit g\"ultige Ausgabe stammt aus dem Jahr 2018.

Die Norm strukturiert den gesamten Produktentwicklungsprozess in einem sogenannten
\emph{V-Modell}: Auf der linken Seite des V erfolgt die Spezifikation von
System-, Hardware- und Software-Anforderungen in absteigender Abstraktionsebene,
auf der rechten Seite die korrespondierende Verifikation und Validierung auf jeder
Ebene \cite{iso26262_2}. Kernkonzept der Norm ist das
\emph{Automotive Safety Integrity Level} (ASIL), eine vierstufige Risikoklassifikation
von ASIL~A (geringste Anforderungen) bis ASIL~D (h\"ochste Anforderungen) sowie die
Klasse QM (Quality Management, keine expliziten Sicherheitsanforderungen).

Die ASIL-Einstufung ergibt sich aus einer Risikoanalyse unter Ber\"ucksichtigung von
Gef\"ahrdungsexposition (E), Kontrollierbarkeit (C) und Schadensausma\ss{} (S). Je
h\"oher der resultierende ASIL, desto strenger die Anforderungen an
Entwicklungsprozess, Verifikation, Testabdeckung und Dokumentation.

\subsection{Anforderungsarten in ISO~26262}
\label{subsec:anforderungsarten}

ISO~26262 unterscheidet verschiedene Klassen von Anforderungen:

\begin{description}[leftmargin=2cm, style=nextline]
  \item[Sicherheitsziele (Safety Goals):] Oberste Anforderungen auf Fahrzeugebene,
    direkt aus der Gefährdungs- und Risikoanalyse (HARA) abgeleitet.
  \item[Funktionale Sicherheitsanforderungen (FSR):] Systemebene, beschreiben das
    geforderte Sicherheitsverhalten.
  \item[Technische Sicherheitsanforderungen (TSR):] Hardware/Software-Ebene,
    beschreiben die technische Implementierung der FSR.
  \item[Software-Sicherheitsanforderungen (SSR):] Spezifisch f\"ur Software-Einheiten,
    direkt implementierbar und formal testbar.
\end{description}

Die Norm verlangt eine l\"uckenlose, bidirektionale
\emph{Anforderungsverfolgbarkeit} (Requirements Traceability) von den Sicherheitszielen
bis hinunter zum einzelnen Softwaremodul und zur\"uck \cite{iso26262_6}.

\subsection{Neuronale Netze -- Grundprinzip und Eigenschaften}
\label{subsec:knn_grundlagen}

K\"unstliche neuronale Netze (KNN) sind mathematische Modelle, die von der Struktur
biologischer Nervensysteme inspiriert sind \cite{Rumelhart1986}. Sie bestehen aus
Schichten von \emph{Neuronen}, deren Verbindungen durch numerische Gewichte
$w_{ij}$ parametriert sind. Die Ausgabe eines Neurons $j$ in Schicht $l$ ergibt sich
aus der gewichteten Summe der Eing\"ange, auf die eine nichtlineare
Aktivierungsfunktion $\sigma$ angewendet wird:

\begin{equation}
  a_j^{(l)} = \sigma\!\left(\sum_{i} w_{ij}^{(l)}\, a_i^{(l-1)} + b_j^{(l)}\right)
  \label{eq:neuron}
\end{equation}

Die Gewichte werden nicht explizit programmiert, sondern durch einen Lernalgorithmus
(typischerweise Backpropagation mit gradientenbasierter Optimierung) aus einem
Trainingsdatensatz \emph{erlernt} \cite{Goodfellow2016}. Dieser Prozess minimiert
eine Verlustfunktion \"uber die Trainingsdaten, ohne dass dem Netz explizites Wissen
\"uber den zugrundeliegenden physikalischen Prozess eingepflanzt wird.

Moderne tiefe Netze (Deep Neural Networks) k\"onnen Millionen bis Milliarden von
Parametern besitzen. Sie haben in zahlreichen Anwendungsgebieten -- Bildklassifikation,
Sprachverarbeitung, Steuerungsaufgaben -- herausragende Leistungen erzielt
\cite{LeCun2015}. Dieser empirische Erfolg sollte jedoch nicht dar\"uber hinwegt\"auschen,
dass die innere Repr\"asentation solcher Netze f\"ur den Menschen im Wesentlichen
undurchschaubar bleibt \cite{Lipton2018}.

\subsection{Der Einsatzkontext: Virtuelle Sensoren auf ECUs}
\label{subsec:virtueller_sensor}

Ein virtueller Sensor ist ein Softwarealgorithmus, der auf Basis priori bekannter
Eingangskan\"ale einen physikalischen Messwert sch\"atzt, der andernfalls nur durch
aufwendige oder kostenintensive Hardware-Sensoren erfasst werden k\"onnte. Die Idee
ist nicht neu -- modellbasierte Beobachter (z.\,B. Luenberger-Beobachter, Kalman-Filter)
sind seit Jahrzehnten im Einsatz \cite{Leveson2011}. Der wesentliche Unterschied
beim KNN-basierten Ansatz: Statt eines physikalisch motivierten, mathematisch
herzuleitenden Modells wird ein rein datengetriebener Ansatz verwendet, dessen
innere Struktur no implizites physikalisches Wissen tr\"agt.

In sicherheitskritischen Anwendungen -- etwa der Sch\"atzung der Fahrzeuggeschwindigkeit,
des Reifenzustands, des Motor- oder Bremsmoments -- kann ein Fehler des virtuellen
Sensors direkt zu Fehlfunktionen sicherheitskritischer Regelsysteme (ABS, ESC,
Lenkhilfe) f\"uhren und damit Leib und Leben gef\"ahrden.

% ============================================================
\section{Die Anforderungen der ISO~26262 im Detail}
\label{sec:iso}
% ============================================================

\subsection{Determinismus}
\label{subsec:determinismus}

ISO~26262 Teil~6 fordert f\"ur ASIL-relevante Software, dass das Softwareverhalten
deterministisch, d.\,h. f\"ur gleiche Eing\"ange jederzeit identisch und vorhersagbar
sein muss \cite{iso26262_6}. Diese Forderung ist keine formale Spitzfindigkeit: Im
Rahmen der Verifikation muss jeder Testfall reproduzierbar sein, Regressionstests
m\"ussen stabile Ergebnisse liefern, und im Falle eines Feldversagens (Field Failure)
muss das Verhalten des Systems auf dem Pr\"ufstand reproduziert werden k\"onnen,
um die Fehlerursache zu identifizieren.

Eine deterministische Ausf\"uhrung setzt voraus, dass:
\begin{itemize}
  \item der Quellcode keine race conditions oder undefiniertes Verhalten enth\"alt,
  \item keine zuf\"alligen Elemente (Zufallszahlen, stochastische Sampling-Verfahren)
    w\"ahrend der Inferenz aktiv sind,
  \item Flie\ss{}kommaoperationen streng nach IEEE~754 und in reproduzierbarer
    Reihenfolge ausgef\"uhrt werden.
\end{itemize}

W\"ahrend ein voll trainiertes, auf festen Gewichten basierendes KNN bei der
Inferenz theoretisch deterministisch ausgef\"uhrt werden kann, ist dies in der
Praxis mit erheblichen H\"urden verbunden. Bereits kleine Unterschiede in der
Reihenfolge von Flie\ss{}kommaoperationen, unterschiedliche Prozessorarchitekturen
(FPU-Implementierungen) oder Compileroptimierungen k\"onnen zu numerisch
unterschiedlichen Ergebnissen f\"uhren -- ein f\"ur die ISO~26262-Verifikation
inakzeptabler Zustand \cite{Dreossi2019}.

\subsection{Anforderungsverfolgbarkeit (Requirements Traceability)}
\label{subsec:traceability}

ISO~26262 Teil~6, Abschnitt~5.4 fordert, dass jede Software-Implementierungseinheit
auf eine oder mehrere Software-Sicherheitsanforderungen zur\"uckf\"uhrbar sein muss
\cite{iso26262_6}. Diese bidirektionale Traceability -- von der Anforderung zum
Code und vom Code zur Anforderung -- ist Voraussetzung f\"ur:

\begin{enumerate}
  \item Die Vollst\"andigkeitsnachweise (sind alle Anforderungen implementiert?),
  \item die Auswirkungsanalyse bei \"Anderungen (Change Impact Analysis),
  \item den Nachweis, dass kein ungeplanter Funktionsumfang (no unintended
    functionality) implementiert wurde.
\end{enumerate}

Ein neuronales Netz besteht nach dem Training aus einer Gewichtsmatrix von
typischerweise Tausenden bis Millionen numerischer Werte. Es ist prinzipiell
nicht m\"oglich, einzelnen Gewichten oder aktivierten Neuronen bestimmte funktionale
Anforderungen zuzuordnen \cite{Lipton2018}. Das Netz als Ganzes \emph{erf\"ullt}
(im Sinne eines black-box-Verhaltens) im g\"unstigsten Fall die Ausgabespezifikation --
die innere Struktur ist jedoch von keiner Anforderung getrieben. Dies stellt eine
fundamentale Verletzung des Traceability-Prinzips dar.

\subsection{Verifikation und Testabdeckung}
\label{subsec:verifikation}

F\"ur ASIL~D-Software schreibt ISO~26262 Teil~6 Tabelle~10 die vollst\"andige
MC/DC-Testabdeckung (Modified Condition/Decision Coverage) vor
\cite{iso26262_6}. MC/DC erfordert, dass jede einzelne Bedingung in einem
zusammengesetzten booleschen Ausdruck den Ausgabewert unabh\"angig beeinflusst.
Diese Anforderung ist auf Quellcode mit expliziter, implementierter Logik ausgerichtet.

Auf ein neuronales Netz l\"asst sich MC/DC nicht anwenden: Es gibt weder
bedingte Verzweigungen im klassischen Sinne noch logische Ausdr\"ucke, deren
Coverage gemessen werden k\"onnte. Alternative Metriken (Neurons-Coverage,
Feature-Space-Coverage) existieren in der Forschung \cite{Huang2017}, sind
jedoch weder in ISO~26262 anerkannt noch in ihrer Aussagekraft bez\"uglich
sicherheitsrelevanter Fehler empirisch belegt.

Dar\"uber hinaus verlangt ISO~26262 die formale Verifikation kritischer
Softwareeigenschaften sowie die vollst\"andige Analyse des Datenflussmodells
(data flow analysis) und Kontrollflusses \cite{iso26262_6}. Beide Analysemethoden
sind f\"ur die Black-Box-Charakteristik neuronaler Netze nicht anwendbar.

\subsection{Keine unbeabsichtigte Funktionalit\"at}
\label{subsec:kein_unbeabsichtigtes}

ISO~26262 verlangt explizit den Nachweis, dass Software keine unbeabsichtigte
Funktionalit\"at (no unintended functionality) implementiert \cite{iso26262_6}.
Diese Anforderung soll sicherstellen, dass das System ausschlie\ss{}lich das tut,
was spezifiziert wurde -- nicht mehr und nicht weniger.

Neuronale Netze lernen ihre internen Repr\"asentationen v\"ollig autonom aus den
Trainingsdaten. Es ist prinzipiell nicht ausgeschlossen, dass das Netz
\emph{Artefakte} oder \emph{Korrelationen} aus den Trainingsdaten internalisiert,
die keine physikalische oder funktionale Grundlage haben, aber das Verhalten
des Netzes f\"ur bestimmte Eingangskonstellationen beeinflussen
\cite{Goodfellow2015}. Da diese internen Repr\"asentationen nicht
einsehbar sind, kann der Nachweis ``keine unbeabsichtigte Funktionalit\"at''
f\"ur ein KNN nicht erbracht werden.

\subsection{\"Anderungsmanagement und Wiederhol-Qualifikation}
\label{subsec:aenderungsmanagement}

ISO~26262 Teil~8 regelt das \"Anderungsmanagement \cite{iso26262_8}. Bei jeder
\"Anderung der Software muss eine \emph{Change Impact Analysis} durchgef\"uhrt
werden, die alle betroffenen Anforderungen und Verifikationsma\ss{}nahmen
identifiziert. Die Bandbreite der notwendigen Requalifikation wird durch die
Auswirkungen der \"Anderung auf das funktionale Sicherheitskonzept bestimmt.

Bei einem KNN bedeutet selbst eine marginale \"Anderung des Trainingsdatensatzes
(z.\,B. Hinzuf\"ugen weniger neuer Messdaten aus dem Feld) eine vollst\"andige
Neu-Parametrierung aller Gewichte \"uber das Gradientenverfahren. Das resultierende
Netz ist ein \emph{neues} Modell, dessen Verhalten systematisch von seinem
Vorg\"anger abweichen kann -- auch in Bereichen, die scheinbar nicht von den
neuen Daten betroffen sind (sog. \emph{catastrophic forgetting} bzw. seine
Umkehrung). Eine klassische, \"anderungsgetriebene Requalifikation ist damit
strukturell nicht m\"oglich; de facto muss bei jeder Modell-Aktualisierung eine
vollst\"andige Neuqualifikation erfolgen.

\subsection{Tool-Qualifikation}
\label{subsec:tool_qualifikation}

ISO~26262 Teil~8, Abschnitt~11 regelt die Qualifikation von Softwareentwicklungswerkzeugen
\cite{iso26262_8}. Der Tool Confidence Level (TCL) bestimmt den erforderlichen Nachweis,
dass ein Werkzeug keine unentdeckten Fehler in die sicherheitsrelevante Software einbringt.

Beim Einsatz von KNNs in der Seriensoftware kommen typischerweise neuartige Trainings-
und Inferenz-Frameworks (TensorFlow, PyTorch, ONNX Runtime) zum Einsatz, die nicht f\"ur
den sicherheitskritischen Einsatz nach ISO~26262 qualifiziert sind. Die TCL-Qualifikation
dieser Werkzeuge ist mit einem erheblichen, kaum bezifferbaren Aufwand verbunden, da die
Werkzeuge selbst komplexe Software mit Millionen von Codezeilen darstellen.

% ============================================================
\section{Spezifische Defizite neuronaler Netze aus Sicherheitsperspektive}
\label{sec:knn_probleme}
% ============================================================

\subsection{Das Erkl\"arbarkeitsdefizit (Explainability Gap)}
\label{subsec:explainability}

Das wohl grunds\"atzlichste Problem neuronaler Netze im Sicherheitskontext ist ihr
\emph{Black-Box-Charakter}: Auch wenn Eingaben und Ausgaben eines Netzes bekannt sind,
bleibt die innere Entscheidungslogik f\"ur den Menschen im Wesentlichen undurchschaubar
\cite{Doshi2017}. Zwar existieren Erkl\"arbarkeitsverfahren wie LIME \cite{Ribeiro2016}
oder SHAP, diese liefern jedoch lediglich \emph{lokale Approximationen} des Netzwerkverhaltens
in der Umgebung eines bestimmten Eingabepunkts -- keine vollst\"andige, formal verifizierbare
Beschreibung des Gesamtverhaltens.

Lipton \cite{Lipton2018} weist darauf hin, dass Interpretierbarkeit kein bin\"ares Merkmal
ist, sondern ein Spektrum mit unterschiedlichen Qualit\"aten. F\"ur den Sicherheitsnachweis
nach ISO~26262 ist jedoch eine \emph{vollst\"andige} Erkl\"arung des Verhaltens
erforderlich -- partielle, statistische Erkl\"arungen reichen nicht aus. Ein
Sicherheitspr\"ufer (Safety Assessor) kann auf Basis von LIME- oder SHAP-Erkl\"arungen
keine sachgem\"a\ss{}e Beurteilung vornehmen.

\subsection{Adversarielle Robustheit}
\label{subsec:adversarial}

Szegedy et al. \cite{Szegedy2014} entdeckten 2014 das Ph\"anomen der
\emph{adversariellen Beispiele}: Minimale, f\"ur das menschliche Auge unsichtbare
St\"orungen eines Eingangsbildes k\"onnen ein tiefes neuronales Netz dazu bringen,
vollst\"andig falsche Vorhersagen mit hoher Konfidenz zu treffen. Goodfellow et al.
\cite{Goodfellow2015} zeigten, dass diese Eigenschaft keine Ausnahme, sondern
ein strukturelles Merkmal von KNNs ist, das sich aus der hochdimensionalen linearen
Natur der erlernten Entscheidungsgrenzen ergibt.

F\"ur sicherheitskritische Sensorapplikationen bedeutet dies: Scheinbar normale,
aber durch Rauschen, Sensor-Drift oder Umgebungsver\"anderungen leicht ver\"anderte
Eingangssignale k\"onnen zu abrupt falschen Sch\"atzwerten f\"uhren -- ohne dass das
System selbst eine erh\"ohte Unsicherheit meldet. Diese Eigenschaft ist aus
sicherheitstechnischer Sicht katastrophal, da sie bedeutet, dass das Versagen des
Systems ohne Vorwarnung und ohne erkennbaren ausl\"osenden Defekt eintreten kann.

\subsection{Out-of-Distribution-Verhalten}
\label{subsec:ood}

Neuronale Netze sind grunds\"atzlich darauf ausgelegt, innerhalb der \emph{Verteilung}
der Trainingsdaten zuverl\"assige Sch\"atzungen zu liefern. Au\ss{}erhalb dieser
Verteilung -- bei Eingaben, die vom Netz noch nicht gesehen wurden (Out-of-Distribution,
OOD) -- ist das Verhalten nicht spezifiziert und h\"aufig fehlerhaft
\cite{Amodei2016}. Dies ist \"aquivalent zum Versagen eines klassischen Software-
moduls bei undefiniertem Eingangs-Wertebereich, mit dem entscheidenden Unterschied:
W\"ahrend ein klassischer Sicherheits-Software-Entwickler explizit Grenzf\"alle
spezifiziert und absichert (\emph{defensive programming}), bieten KNNs keinerlei
Mechanismus, OOD-Situationen zu erkennen oder abzusichern -- au\ss{}er durch
aufwendige, nachtr\"aglich hinzugef\"ugte OOD-Erkennungsmodule.

Koopman und Fratrik \cite{Koopman2019} zeigen, dass die praktisch unendliche
Vielfalt von OOD-Szenarien in der realen Welt -- ver\"anderte Wetterbedingungen,
unbekannte Stra\ss{}enbedingungen, unerwartete Objektkonstellationen -- eine
vollst\"andige Testabdeckung des Eingangsraums prinzipiell ausschlie\ss{}t.

\subsection{Fehlende formale Spezifikation}
\label{subsec:formale_spez}

Klassische sicherheitskritische Software wird gegen eine \emph{formale Spezifikation}
entwickelt: Das Verhalten der Software ist f\"ur alle relevanten Eingangs-
kombinationen eindeutig definiert \cite{Seshia2018}. Dieses Grundprinzip ist in
ISO~26262 fest verankert und bildet die Basis f\"ur jeden Verifikationsnachweis.

Bei einem KNN existiert keine solche formale Spezifikation: Das Training
maximiert eine Verlustfunktion \"uber einen endlichen Datensatz und erzeugt
damit ein Modell, dessen Verhalten au\ss{}erhalb des Trainingsdatensatzes
lediglich durch Generalisierungsannahmen und empirische Tests charakterisiert
werden kann. Ein mathematisch vollst\"andiger Nachweis der Korrektheit
f\"ur den gesamten Eingangsraum ist -- abgesehen von kleinen, niedrigdimensionalen
Netzen -- mit dem Stand der Technik nicht m\"oglich \cite{Huang2017}.

\subsection{Statistische vs.\ deterministische Korrektheit}
\label{subsec:stat_det}

Ein zentrales Missverst\"andnis beim Einsatz von KNNs in
sicherheitskritischen Systemen ist die implizite Gleichsetzung von
\emph{statistischer Genauigkeit} und \emph{sicherheitstechnischer Korrektheit}.
Ein KNN, das in 99{,}9\% aller Testf\"alle korrekte Ergebnisse liefert, zeigt
im verbleibenden 0{,}1\% ein undefiniertes Verhalten -- und k\"onnte gerade in
diesen Randf\"allen in den kritischsten Fahrsituationen versagen. Die Norm
ISO~26262 kennt keine statistischen Fehlerquoten f\"ur Softwarefehler (im
Unterschied zu Hardware-Ausfallraten); Softwarekorrektheit wird als Ja/Nein-Eigenschaft
betrachtet, die durch formale Verifikation und Testabdeckung nachzuweisen ist
\cite{iso26262_6}.

Varshney und Alemzadeh \cite{Varshney2017} verdeutlichen diesen Unterschied:
\"In safety-critical machine learning applications, one misbehavior may be
far more costly than one thousand correct behaviors.`` Diese Asymmetrie ist im
Sicherheitsdenken von ISO~26262 fundamental -- und mit dem statistischen
Qualit\"atsbegriff von KNNs unvereinbar.

\subsection{Reproduzierbares Schadensversagen (Fault Reproduction)}
\label{subsec:fault_repro}

Ein wesentliches Element der funktionalen Sicherheitsentwicklung ist die F\"ahigkeit,
im Feldversagensfall das fehlerhafte Verhalten reproduzieren und analysieren zu
k\"onnen (\emph{fault reproduction}), um die Ursache zu identifizieren und durch
Softwarekorrektur zu beseitigen. Bei klassischer Software ist dies m\"oglich, da
der Code deterministisch und vollst\"andig \"uberpr\"ufbar ist.

Bei einem KNN-basierten virtuellen Sensor im Feld kann ein Versagen durch nichtlineare
Wechselwirkungen multidimensionaler Eingangssignale verursacht sein, die im
Laborumfeld nicht reproduzierbar sind. Selbst wenn der Eingangsdatenvektor des
Versagenszeitpunkts vollst\"andig rekonstruiert werden kann, ist eine kausale Analyse
-- warum hat das Netz genau diesen falschen Wert ausgegeben? -- nicht m\"oglich.
Dies untergr\"abt ein Kernprinzip des Sicherheitsprozesses, n\"amlich das
kontinuierliche Lernen aus Feldversagensf\"allen zur Verbesserung der Systemsicherheit
\cite{Leveson2011}.

\subsection{Zertifizierungsnachweis und Safety Case}
\label{subsec:safety_case}

ISO~26262 Teil~2 verlangt die Erstellung eines \emph{Safety Cases} -- eines
strukturierten Arguments, das auf Basis von Evidenz nachweist, dass das System
ausreichend sicher ist \cite{iso26262_2}. Dieser Safety Case muss von einem
unabh\"angigen Sicherheitsprüfer (Functional Safety Assessor) bewertet und
best\"atigt werden.

F\"ur ein KNN-basiertes System l\"asst sich kein vollst\"andiger Safety Case
erstellen: Die notwendige Evidenz -- insbesondere der Nachweis der formalen
Verifikation und der vollst\"andigen Anforderungsabdeckung -- kann f\"ur KNNs
strukturell nicht erbracht werden \cite{AMLAS2021}. Bestehende Rahmenwerke wie
AMLAS (Assurance of Machine Learning in Autonomous Systems) \cite{AMLAS2021}
stellen fest, dass f\"ur ML-basierte Systeme neue, noch nicht standardisierte
Nachweisformen erforderlich sind -- womit sie implizit einr\"aumen, dass mit den
Mitteln der ISO~26262 allein kein vollst\"andiger Nachweis m\"oglich ist.

% ============================================================
\section{Zentrale Argumente gegen den Einsatz}
\label{sec:argumente}
% ============================================================

Aus den vorstehenden Analysen lassen sich die folgenden zentralen Argumente
destillieren, die konsequent gegen den Einsatz von neuronalen Netzen in
ISO~26262-konformer, sicherheitskritischer ECU-Software sprechen:

\subsection{Argument 1: Strukturelle Unvereinbarkeit mit dem Verifikationsparadigma}
\label{subsec:arg1}

ISO~26262 basiert auf einem \emph{konstruktiven Verifikationsparadigma}: Software
wird anhand expliziter Anforderungen entwickelt und formal gegen diese Anforderungen
verifiziert. Neuronale Netze sind das Produkt eines \emph{empirischen Optimierungsparadigmas}:
Sie entstehen durch Minimierung einer Verlustfunktion \"uber einen endlichen Datensatz.
Diese beiden Paradigmen sind fundamental unvertr\"aglich. Es handelt sich nicht um eine
L\"ucke, die durch h\"oheren Aufwand oder bessere Werkzeuge geschlossen werden kann,
sondern um eine epistemologische Grundverschiedenheit: Sicherheitsnachweis beruht
auf \emph{Wissen}, KNN-Training erzeugt \emph{Empirie}. Wissen ist übertragbar,
begr\"undbar und vollst\"andig; Empirie ist stets begrenzt, stichprobenabh\"angig
und niemals vollst\"andig.

\subsection{Argument 2: Vertrauensw\"urdigkeit erfordert Erkl\"arbarkeit}
\label{subsec:arg2}

Vertrauensw\"urdigkeit (\emph{trustworthiness}) ist nach Avizienis et al.
\cite{Avizienis2004} eine zusammengesetzte Eigenschaft, die Zuverl\"assigkeit
(\emph{reliability}), Verf\"ugbarkeit (\emph{availability}), Wartbarkeit
(\emph{maintainability}) und Sicherheit (\emph{safety}) umfasst. F\"ur sicherheitskritische
Systeme setzt Vertrauensw\"urdigkeit voraus, dass das System nicht nur in der
Vergangenheit korrekt funktioniert hat, sondern dass \emph{begr\"undet werden kann},
warum es in Zukunft korrekt funktionieren wird -- insbesondere auch in Situationen,
die noch nicht getestet wurden.

Neuronale Netze k\"onnen diesen Nachweis nicht liefern: Ihr Erfolg in
Testszenarien belegt die vergangene Leistung, nicht die k\"unftige Korrekheit
in unbekannten Situationen. Wing \cite{Wing2021} formuliert f\"ur vertrauensw\"urdige
KI-Systeme explizit die Anforderungen \emph{fair, robust, explainable, private,
safe and secure} -- und stellt fest, dass kein heutiges KNN diese Anforderungen
vollst\"andig erf\"ullt. Im Kontext sicherheitskritischer Systeme ist fehlende
Vertrauensw\"urdigkeit kein theoretisches Problem, sondern ein rechtlich relevanter
Haftungsausschlussgrund.

\subsection{Argument 3: Determinismusdefizienz als dauerhaftes Strukturproblem}
\label{subsec:arg3}

Der Determinismusbegriff in ISO~26262 ist nicht auf das laufzeitliche
Ausf\"uhrungsverhalten beschr\"ankt, sondern umfasst auch die
\emph{Vorhersagbarkeit des Verhaltens \"uber den gesamten Betriebsbereich}. Ein
System, dessen Verhalten vollst\"andig bekannt und berechenbar ist, ist deterministisch;
ein System, dessen Verhalten f\"ur bestimmte Eingaben unbekannt oder unvorhersehbar ist,
ist es nicht. Da f\"ur KNNs das Verhalten au\ss{}erhalb der Trainingsverteilung
undefiniert ist und das exakte Verhalten f\"ur beliebige Eingaben nicht
vorausberechnet werden kann, sind KNNs im sicherheitstechnisch relevanten Sinne
nicht deterministisch \cite{Dreossi2019}.

\subsection{Argument 4: Datenqualit\"at als neue, unkontrollierbare Fehlerquelle}
\label{subsec:arg4}

Klassische sicherheitskritische Software enth\"alt keine externe Abh\"angigkeit
zu einem Datensatz w\"ahrend der Funktionsausf\"uhrung. Das Verhalten der
Software ist vollst\"andig durch den Quellcode bestimmt, der versioniert, auditiert
und in seiner Vollst\"andigkeit den Entwicklernormen entsprechend spezifiziert
ist \cite{Sommerville2016}.

KNN-Modelle hingegen repr\"asentieren das implizite Wissen des Trainingsdatensatzes.
Die Qualit\"at dieses Datensatzes -- seine Vollst\"andigkeit, seine Repr\"asentativit\"at,
seine Fehlerfreiheit -- hat direkten Einfluss auf das Verhalten des trainierten Netzes
im Feld. ISO~26262 enth\"alt keine Methodik f\"ur die Qualit\"atssicherung von
Trainingsdaten, da dieses Konzept im Entwicklungszeitraum der Norm nicht relevant war.
Damit entsteht eine neue Klasse von Fehlerquellen -- \emph{Datenfehler} -- die
au\ss{}erhalb des normativen Rahmens liegen und damit unkontrolliert bleiben
\cite{Gauerhof2018}. Fehlerhafte, veraltete oder nicht-repr\"asentative Trainingsdaten
k\"onnen systematic errors im Modell erzeugen, die mit keinem der in ISO~26262
vorgeschriebenen Verifikationsverfahren aufgedeckt werden k\"onnen.

\subsection{Argument 5: Systemische \"Uberschaubarkeit}
\label{subsec:arg5}

ISO~26262 verlangt nicht nur die Sicherheit einzelner Komponenten, sondern die
\"Uberschaubarkeit des Gesamtsystems (\emph{system comprehensibility}), die
Voraussetzung f\"ur eine korrekte und vollst\"andige Sicherheitsanalyse (FMEA,
FTA) ist \cite{iso26262_9}. Ein System, das ein KNN enth\"alt, dessen
innere Funktionsweise nicht vollst\"andig durchleuchtet werden kann, ist per
definitionem nicht vollst\"andig \"uberschaubar.

Konsequenz: Eine FMEA (Fehlerm\"oglichkeits- und Einflussanalyse) f\"ur ein KNN-Modul
kann keine vollst\"andige Fehlerliste erstellen, da die Fehlermodi eines Netzes --
welche Eingangskonstellationen f\"uhren zu welchen Fehlausgaben? -- nicht a priori
spezifizierbar sind. Damit ist die Vollst\"andigkeit jeder Sicherheitsanalyse,
die KNN-Komponenten enth\"alt, grundlegend in Frage gestellt.

\subsection{Argument 6: Haftungs- und Nachweisrechtliche Implikationen}
\label{subsec:arg6}

In Deutschland und der EU unterstehen Fahrzeuge dem Produkthaftungsrecht
(ProdHaftG, EU-Produkthaftungsrichtlinie 2024/2853). Im Falle eines
Schadensersatzanspruchs muss der Hersteller nachweisen, dass das Fahrzeug
\emph{alle zum Zeitpunkt des Inverkehrbringens relevanten Sicherheitsstandards
erf\"ullt hat}. ISO~26262 ist der ma\ss{}gebliche Industriestandard und wird
von Typgenehmigungsbeh\"orden als Referenz herangezogen.

Kann ein Hersteller die Normkonformit\"at seines KNN-basierten virtuellen Sensors
nicht l\"uckenlos nachweisen -- und dieser Nachweis ist, wie in den vorstehenden
Abschnitten gezeigt, nur partiell m\"oglich --, so tr\"agt er ein erhebliches
Haftungsrisiko. Im schlimmsten Fall -- einem durch den KNN-Ausfall verursachten
Personenschaden -- ist der fehlende Sicherheitsnachweis als Indiz f\"ur fahrl\"assige
Konstruktion zu werten. Dieses Risiko ist in der \"offentlichen und industriellen
Diskussion bislang unterbewertet.

\subsection{Argument 7: Zertifizierungsrisiko bei Typgenehmigung}
\label{subsec:arg7}

Fahrzeuge, die auf europa\"aischen Stra\ss{}en betrieben werden, ben\"otigen eine
Typgenehmigung, die zunehmend die Einhaltung funktionaler Sicherheitsstandards
voraussetzt (UNECE WP.29, ALKS-Verordnung). Typgenehmigungsbeh\"orden wie das
Kraftfahrtbundesamt (KBA) oder die britische DVSA bewerten zunehmend KI-basierte
Systeme, verf\"ugen jedoch noch \"uber keine konsolidierten Leitf\"aden f\"ur deren
Bewertung. Die resultierende Unsicherheit bedeutet: Ein Fahrzeughersteller,
der KNNs in ASIL-relevanten Systemen einsetzt, muss mit l\"angeren, teureren und
unsichereren Typgenehmigungsverfahren rechnen -- und l\"auft Gefahr, nach
Markteinf\"uhrung mit Nachr\"ustungsanforderungen konfrontiert zu werden, wenn
die Beh\"orden ihre Bewertungspraxis konsolidieren.

% ============================================================
\section{Kritische Auseinandersetzung}
\label{sec:gegenargumente}
% ============================================================

\subsection{ASIL-Dekomposition und Monitoring-Architektur}
\label{subsec:gegenarg1}

Das am h\"aufigsten vorgebrachte Argument f\"ur die Verwendbarkeit von KNNs in
sicherheitskritischen Systemen ist die \emph{ASIL-Dekomposition}: Das KNN wird auf
QM oder ASIL~B eingestuft und durch einen deterministischen, ASIL~D-qualifizierten
Monitor \"uberwacht, der die Ausgabe des KNN auf Plausibilit\"at pr\"uft und im
Fehlerfall auf einen sicheren Zustand umschaltet.

Dieser Ansatz ist prinzipiell legitim und in ISO~26262 vorgesehen
\cite{iso26262_9}. Er hat jedoch entscheidende Grenzen:

\begin{enumerate}
  \item \textbf{Monitor-Vollst\"andigkeit:} Der Monitor kann nur Fehler erkennen,
    die er zu erkennen in der Lage ist. F\"ur die Spezifikation des Monitors m\"ussen
    die m\"oglichen Fehlm\"odi des KNN bekannt sein -- was, wie in
    \cref{subsec:arg5} gezeigt, gerade nicht der Fall ist. Ein Monitor, dessen
    Fehlermodell unvollst\"andig ist, bietet keine vollst\"andige Sicherheitsgarantie.

  \item \textbf{Gemeinsam-Ursachen-Ausfall (CCF):} Wenn KNN und Monitor beide
    auf denselben Eingangsdaten arbeiten und das KNN durch adversarielle oder
    OOD-Eingaben gest\"ort wird, k\"onnen dieselben Eingaben auch den Monitor
    t\"auschen, sofern dieser kein v\"ollig unabh\"angiges Funktionsprinzip verwendet.

  \item \textbf{Funktionale Degradation:} Ein Monitoring-basiertes System degradiert
    den KNN-Ansatz faktisch zum Hilfssystem mit eingeschr\"anktem ASIL. Wenn die
    eigentliche Sicherheit ausschlie\ss{}lich durch den Monitor gew\"ahrleistet wird,
    stellt sich die Frage, welchen Merkmalsvorteil das KNN gegen\"uber einem
    klassischen, deterministischen Sch\"atzer noch bietet.
\end{enumerate}

\subsection{ISO/PAS~8800 und neue Normen adressieren das Problem}
\label{subsec:gegenarg2}

Die 2024 ver\"offentlichte ISO/PAS~8800 \emph{Road Vehicles -- Safety and Artificial
Intelligence} \cite{isopas8800} wird als normative L\"osung f\"ur die Integration von
KI in sicherheitskritische Fahrzeugsysteme beworben. In der Tat adressiert ISO/PAS~8800
viele der in dieser Arbeit genannten Probleme und definiert ML-spezifische Work Products,
Datenqualit\"atsanforderungen und ODD-Spezifikationen.

Allerdings ist kritisch anzumerken:

\begin{enumerate}
  \item ISO/PAS~8800 ist ein \emph{Publicly Available Specification} (PAS) -- eine
    vorl\"aufige Norm, kein vollwertiger ISO-Standard. Sie hat noch nicht den
    Reifegrad einer mehrfach durchlaufenen Konsensent\"scheidung.
  \item Die Norm definiert Prozessanforderungen, l\"ost jedoch nicht die
    \emph{technischen} Grundprobleme -- sie spezifiziert lediglich, wie der
    Nachweis organisiert werden soll. Die Frage, ob ein statistisches Modell
    wirklich sicher agiert, wird nicht durch bessere Prozessorganisation beantwortet.
  \item ISO/PAS~8800 muss im Zusammenspiel mit ISO~26262 gelesen werden und
    \emph{ersetzt} diese nicht. Soweit ISO~26262 Anforderungen stellt, die
    KNNs strukturell nicht erf\"ullen k\"onnen, \"andert ISO/PAS~8800 daran nichts.
\end{enumerate}

\subsection{MIL/SIL/PIL-Tests belegen die Sicherheit}
\label{subsec:gegenarg3}

Toolanbieter wie dSPACE (TargetLink) betonen die MIL/SIL/PIL-Testkette als Nachweis,
dass das KNN auf der Zielhardware genauso funktioniert wie im Modell. Dies ist ein
wichtiger, aber h\"aufig missdeuteter Beitrag:

MIL/SIL/PIL-Tests belegen \emph{Implementierungskonsistenz} -- das hei\ss{}t,
das trainierte Modell wird durch den Codegenerierungsprozess nicht verf\"alscht.
Sie sagen jedoch \emph{nichts} \"uber die Korrektheit des trainierten Modells
selbst aus. Das korrekte Deployment eines fehlerhaften Modells ist kein
Sicherheitsgewinn.

\subsection{Akzeptiertes Restrisiko (Acceptable Residual Risk)}
\label{subsec:gegenarg4}

Ein pragmatisches Argument lautet: ISO~26262 verlangt nicht null Fehler, sondern
nur eine Reduzierung des Risikos auf ein akzeptables Ma\ss{} (\emph{ALARP:
As Low As Reasonably Practicable}). Neuronale Netze als solche seien akzeptabel,
solange das Gesamtrisiko unterhalb des Schwellenwerts bleibt.

Dieses Argument verkennt, dass ALARP unter ISO~26262 keine statistische Abw\"agung
von Fehlerraten ist, sondern einen \emph{Entwicklungsprozessnachweis} erfordert:
Das Risiko gilt als ausreichend reduziert, wenn alle normativen Anforderungen an
den Entwicklungsprozess erf\"ullt wurden. Da diese Anforderungen f\"ur KNNs nicht
vollst\"andig erf\"ullt werden k\"onnen, kann ALARP nicht beansprucht werden --
unabh\"angig von der empirisch gemessenen Fehlerrate des Netzes.

% ============================================================
\section{Sichere Alternativen zu KNN-basierten Sensorl\"osungen}
\label{sec:alternativen}
% ============================================================

Die Ablehnung von KNNs f\"ur sicherheitskritische Funktionen schlie\ss{}t den
Einsatz moderner, leistungsf\"ahiger Technik nicht aus. Folgende Alternativen sind
mit ISO~26262 vereinbar und bieten f\"ur viele Anwendungsbereiche gleichwertige
oder \"uberlegene L\"osungen:

\subsection{Modellbasierte Beobachter und Kalman-Filter}
\label{subsec:alt_kalman}

F\"ur die Sch\"atzung physikalischer Gr\"o\ss{}en aus Messdaten sind
\emph{Kalman-Filter} und deren nichtlineare Erweiterungen (Extended Kalman Filter,
Unscented Kalman Filter) seit Jahrzehnten der Stand der Technik in
sicherheitskritischen Systemen \cite{Leveson2011}. Sie basieren auf einem expliziten
physikalischen Systemmodell, sind linear in den Zustandsgr\"o\ss{}en (bzw. lokal
linearisiert), haben eine bekannte, analytisch berechenbare Fehlerkovarianz und
erm\"oglichen eine vollst\"andige formale Analyse. Die Anforderungsverfolgbarkeit
ist l\"uckenlos: jeder Parameter des Filters leitet sich aus dem zugrundeliegenden
physikalischen Modell ab.

\subsection{Polynomiale Approximationen und Look-up-Tabellen}
\label{subsec:alt_lut}

F\"ur nichtlineare Kennlinien und Kennfelder -- h\"aufige Anwendungsf\"alle f\"ur
virtuelle Sensoren -- sind stufenweise polynomiale Approximationen und
Look-up-Tabellen mit Interpolation eine bew\"ahrte, sicherheitstechnisch
unproblematische L\"osung. Das Verhalten ist vollst\"andig durch explizite Werte
definiert, \"uberpr\"ufbar und formal verifizierbar.

\subsection{Physikalisch-motivierte Modelle}
\label{subsec:alt_phys}

In vielen F\"allen, in denen ein neuronales Netz als virtueller Sensor vorgeschlagen
wird, liegt hinreichendes physikalisches Wissen vor, um ein parametrisches Modell
abzuleiten. Der Einsatz von KNNs wird dann h\"aufig mit geringerem Modellierungsaufwand
begr\"undet. Dieser Aufwand ist jedoch eine Investition in Sicherheit und
Nachvollziehbarkeit -- er ist im Sicherheitsentwicklungsprozess nicht wegoptimierbar.

\subsection{Fuzzy-Logik und Expertensysteme}
\label{subsec:alt_fuzzy}

F\"ur qualitative, wissensbasierte Entscheidungsfunktionen boten Fuzzy-Logik-Systeme
und Expertensysteme schon l\"ange erkl\"arbare, formal analysierbare Alternativen
zu KNNs, die mit ISO~26262 vereinbar sind \cite{Kurd2003}. Sie sind in
eingebetteten Systemen effizient implementierbar und bieten eine h\"ohere
Erkl\"arbarkeit.

\subsection{Formale Methoden und Model Checking}
\label{subsec:alt_formal}

F\"ur die Verifikation kritischer Systemeigenschaften -- Sicherheitsinvarianten,
zeitliche Eigenschaften -- stehen formale Methoden (Model Checking, Theorem
Proving) zur Verf\"ugung, die in ISO~26262 ASIL~D explizit empfohlen werden
\cite{iso26262_6}. Diese Methoden sind auf klassisch programmierten Code gut
anwendbar und liefern mathematisch gesicherte Nachweise.

% ============================================================
\section{Diskussion: Kumuliertes Defizit und Normative L\"ucke}
\label{sec:diskussion}
% ============================================================

Die in dieser Arbeit analysierten Probleme sind keine isolierten technischen
Einzelprobleme, die durch spezifische Abhilfe\-ma\ss{}nahmen behebbar w\"aren.
Sie repr\"asentieren ein \emph{kumuliertes, strukturelles Defizit}, das aus der
fundamentalen Unvereinbarkeit zweier Designphilosophien resultiert:

Die sicherheitstechnische Entwicklungsphilosophie der ISO~26262 beruht auf dem
Prinzip: \emph{Alles, was nicht explizit spezifiziert und verifiziert wurde,
ist als unsicher zu betrachten.} Die maschinelle Lernphilosophie neuronaler Netze beruht auf dem Prinzip: \emph{Was empirisch gut genug funktioniert, kann eingesetzt werden.} Diese beiden Prinzipien sind nicht kompatibel. Sie k\"onnen durch keinen technischen Trick oder normativen Workaround in \"Ubereinstimmung gebracht werden.

\subsection{Die normative L\"ucke}
\label{subsec:normative_luecke}

Die Gemeinschaft der Sicherheitsnormierer erkennt dieses Dilemma zunehmend an.
ISO~21448 (SOTIF) \cite{iso21448} adressiert Versagen durch unzureichende
Spezifikation oder Leistung -- deckt aber nicht die strukturellen
Verifikationsdefizite von KNNs ab. ISO/PAS~8800 \cite{isopas8800} definiert
ML-spezifische Prozessanforderungen, ersetzt jedoch nicht die formalen
Verifikationsanforderungen der ISO~26262. AMLAS \cite{AMLAS2021} bietet
einen Argumentationsrahmen f\"ur ML-Safety-Cases, r\"aumt aber selbst ein,
dass wesentliche Nachweisformen noch nicht standardisiert sind.

In Summe besteht eine \emph{normative L\"ucke}: Es gibt keinen anerkannten, vollst\"andigen
Sicherheitsnachweis f\"ur KNN-Komponenten in ASIL-relevanten Steuerger\"atefunktionen.
Diese L\"ucke schlie\ss{}t Werkzeugketten (TargetLink, ONNX-Deployment) nicht:
Das korrekte Deployment eines Modells bedeutet nicht dessen Sicherheitsnachweis.

\subsection{Das Pr\"azedenzargument der Luft- und Raumfahrt}
\label{subsec:praezedenz}

Interessant ist der Blick auf die Luft- und Raumfahrt, wo funktionale Sicherheit
\"ahnlich streng gehandhabt wird wie im Automobil (DO-178C als Pendant zu ISO~26262).
Das Neural Network Certification Advisory Circular der FAA (FAA-AC) hat explizit
festgestellt, dass neuronale Netze f\"ur sicherheitskritische Flugsteuerungssysteme
nicht zertifizierbar sind und begr\"undet dies mit denselben Argumenten wie diese
Arbeit: fehlende Erkl\"arbarkeit, nicht anwendbare Coverage-Metriken und
unvollst\"andige Fehlermodelle. Erst in j\"ungster Zeit (2023-2024) werden erste
Pilotprojekte unter eng umrissenen Einschr\"ankungen und mit aufwendigen
Kompensationsma\ss{}nahmen zugelassen -- und dies nach Jahren intensiver
Standardisierungsarbeit, die im Automobil noch nicht abgeschlossen ist.

\subsection{Das Tempo des technologischen Wandels vs.\ normative Reife}
\label{subsec:tempo}

Ein oft unterschlager Aspekt der Diskussion ist die zeitliche Dimension: KI und
Deep Learning entwickeln sich in einem Tempo, das weit \"uber die Entwicklungszyklen
von Sicherheitsnormen hinausgeht. ISO~26262:2018 wurde nach einem Jahrzehnt
internationaler Konsensfindung ver\"offentlicht. ISO/PAS~8800:2024 ist noch nicht
in die ISO~26262 integriert. Bis ein konsolidiertes, internationales Regelwerk
f\"ur KI in sicherheitskritischen Automobil\-systemen existiert, das alle hier
diskutierten Defizite adressiert, werden weitere Jahre vergehen -- Jahre, in denen
Fahrzeuge mit KNN-basierten Funktionen bereits auf den Stra\ss{}en fahren k\"onnten.

Die Verantwortung, in dieser Unsicherheitsphase konservativ zu handeln -- also
KNNs aus sicherheitskritischen Funktionen fernzuhalten oder auf gut begr\"undete,
eng beschr\"ankte Ausnahmen zu reduzieren --, liegt bei den Fahrzeugherstellern
und Tier-1-Zulieferern. Das Argument des Wettbewerbsdrucks ist dabei kein
Sicherheitsargument.

% ============================================================
\section{Fazit und Empfehlungen}
\label{sec:fazit}
% ============================================================

Die vorliegende Arbeit hat eine systematische Analyse der Vereinbarkeit von
neuronalen Netzen mit den Anforderungen des Sicherheitsstandards ISO~26262
vorgenommen. Das Ergebnis ist eindeutig: Neuronale Netze in ihrer gegenw\"artigen
Form erf\"ullen die Kernanforderungen der ISO~26262 an ASIL-relevante Software
strukturell nicht. Die identifizierten Defizite betreffen:

\begin{itemize}[itemsep=2pt]
  \item fehlenden Determinismus im sicherheitstechnisch relevanten Sinne,
  \item unm\"ogliche Anforderungsverfolgbarkeit innerhalb der Netzwerkstruktur,
  \item Nicht-Anwendbarkeit der vorgeschriebenen Verifikationsverfahren (MC/DC),
  \item unbeweisbare Abwesenheit unbeabsichtigter Funktionalit\"at,
  \item strukturell unkontrolliertes OOD- und adversarielle Verhalten,
  \item unvollst\"andige Fehlermodelle und damit unvollst\"andige Sicherheitsanalysen,
  \item fehlende normative und werkzeugm\"a\ss{}ige Grundlage f\"ur vollst\"andige
    Safety Cases.
\end{itemize}

Diese Defizite k\"onnen durch systemarchitektonische Kompensationsma\ss{}nahmen
(Monitoring, ASIL-Dekomposition) teilweise abgemildert werden, verschwinden jedoch
nicht. Die \"ubergeordnete epistemische Grenze bleibt: Ein statistisch erlerntes
Modell kann nie mit derselben Sicherheitsevidenz versehen werden wie ein formal
spezifiziertes, deterministisch implementiertes und vollst\"andig verifiziertes
Softwaremodul.

\bigskip
\noindent\textbf{Empfehlungen:}

\begin{enumerate}[itemsep=4pt]
  \item \textbf{Keine KNN-Nutzung in ASIL~C/D-Funktionen:} F\"ur sicherheitskritische
    Funktionen mit ASIL~C- oder ASIL~D-Einstufung sollten ausschlie\ss{}lich klassische,
    formal verifizierbare Algorithmen eingesetzt werden, bis ein verabschiedeter,
    vollst\"andiger normativer Rahmen f\"ur KI in sicherheitskritischen Systemen vorliegt.

  \item \textbf{Beschr\"ankung auf QM / unkritische Hilfsf\"unktionen:} KNNs k\"onnen
    sinnvoll in der Vorverarbeitung von Rohdaten, in der Kalibrierung von
    Diagnosealgorithen oder in der ASIL-QM-Komfortfunktionalit\"at eingesetzt werden,
    wo keine direkten Sicherheitsziele betroffen sind.

  \item \textbf{Normative Grundlage abwarten und mitgestalten:} Fahrzeughersteller
    und Zulieferer sollten aktiv an der Weiterentwicklung von ISO/PAS~8800,
    ISO~21448 und verwandten Normen mitwirken, anstatt bestehende normative L\"ucken
    als implizite Erlaubnis zu interpretieren.

  \item \textbf{Investition in erkl\"arbare und verifizierbare KI (XAI):} Die
    Forschung zu formal verifizierbaren, erkl\"arbaren KI-Systemen (z.\,B.
    neurosymbolische Systeme, formale Verifikation kleiner Netze) sollte
    intensiviert und mit der Normarbeit verzahnt werden, bevor ein breiter
    Einsatz in ASIL-relevanten Funktionen erfolgt.

  \item \textbf{Werkzeugkettenqualifikation:} Wer KNNs trotz der genannten
    Einschr\"ankungen einsetzt, muss die vollst\"andige TCL-Qualifikation der
    verwendeten KI-Frameworks sicherstellen -- ein Aufwand, der in der Praxis
    h\"aufig unterschätzt wird und in die Kosten-Nutzen-Bewertung einzubeziehen ist.
\end{enumerate}

\bigskip
\noindent Diese Arbeit tritt nicht f\"ur eine pauschale Ablehnung von KI im
Fahrzeug ein. In zahlreichen Anwendungsgebieten -- Komfortfunktionen, Infotainment,
lernende Fahrerassistenz ohne direkten Sicherheits\-bezug -- kann und soll KI ihr
Potenzial entfalten. F\"ur sicherheitskritische Steuerger\"atef\"unktionen aber
gilt: Die Last des Sicherheitsnachweises liegt beim Entwickler, nicht beim Pr\"ufer.
Solange dieser Nachweis f\"ur neuronale Netze nicht vollst\"andig und konform zu
ISO~26262 erbracht werden kann, verbietet sich ihr Einsatz in solchen Funktionen
aus Gr\"unden der technischen Vernunft, der Rechtssicherheit und -- zuvorderst --
des Schutzes von Menschenleben.

\begin{thebibliography}{99}
\addcontentsline{toc}{section}{Literaturverzeichnis}

\bibitem{iso26262_1}
International Organization for Standardization (ISO):
\textit{ISO~26262-1: Road vehicles -- Functional safety -- Part~1: Vocabulary}.
2.~Ausgabe. Genf: ISO, 2018.

\bibitem{iso26262_2}
International Organization for Standardization (ISO):
\textit{ISO~26262-2: Road vehicles -- Functional safety -- Part~2: Management of functional safety}.
2.~Ausgabe. Genf: ISO, 2018.

\bibitem{iso26262_4}
International Organization for Standardization (ISO):
\textit{ISO~26262-4: Road vehicles -- Functional safety -- Part~4: Product development at the system level}.
2.~Ausgabe. Genf: ISO, 2018.

\bibitem{iso26262_6}
International Organization for Standardization (ISO):
\textit{ISO~26262-6: Road vehicles -- Functional safety -- Part~6: Product development at the software level}.
2.~Ausgabe. Genf: ISO, 2018.

\bibitem{iso26262_8}
International Organization for Standardization (ISO):
\textit{ISO~26262-8: Road vehicles -- Functional safety -- Part~8: Supporting processes}.
2.~Ausgabe. Genf: ISO, 2018.

\bibitem{iso26262_9}
International Organization for Standardization (ISO):
\textit{ISO~26262-9: Road vehicles -- Functional safety -- Part~9: ASIL-oriented and safety-oriented analyses}.
2.~Ausgabe. Genf: ISO, 2018.

\bibitem{iso21448}
International Organization for Standardization (ISO):
\textit{ISO~21448: Road vehicles -- Safety of the intended functionality (SOTIF)}.
Genf: ISO, 2022.

\bibitem{isopas8800}
International Organization for Standardization (ISO):
\textit{ISO/PAS~8800: Road vehicles -- Safety and artificial intelligence}.
Genf: ISO, 2024.

\bibitem{iec61508}
International Electrotechnical Commission (IEC):
\textit{IEC~61508: Functional Safety of Electrical/Electronic/Programmable Electronic Safety-related Systems}.
2.~Ausgabe. Genf: IEC, 2010.

\bibitem{AMLAS2021}
Hawkins, R.; Paterson, C.; Picardi, C.; Jia, Y.; Calinescu, R.; Habli, I.:
\textit{AMLAS: A Methodology for the Assurance of Machine Learning in Autonomous Systems}.
University of York, 2021.
arXiv:2102.10020.

\bibitem{Salay2018}
Salay, R.; Khastgir, S.; Czarnecki, K.:
An Analysis of ISO~26262: Using Machine Learning Safely in Automotive Software.
\textit{arXiv preprint arXiv:1709.02435}, 2018.

\bibitem{Koopman2016}
Koopman, P.; Wagner, M.:
Challenges in Autonomous Vehicle Testing and Validation.
In: \textit{SAE World Congress}. SAE International, 2016.

\bibitem{Amodei2016}
Amodei, D.; Olah, C.; Steinhardt, J.; Christiano, P.; Schulman, J.; Man\'e, D.:
Concrete Problems in AI Safety.
\textit{arXiv preprint arXiv:1606.06565}, 2016.

\bibitem{Goodfellow2015}
Goodfellow, I.; Shlens, J.; Szegedy, C.:
Explaining and Harnessing Adversarial Examples.
In: \textit{International Conference on Learning Representations (ICLR)}, 2015.

\bibitem{Szegedy2014}
Szegedy, C.; Zaremba, W.; Sutskever, I.; Bruna, J.; Erhan, D.; Goodfellow, I.; Fergus, R.:
Intriguing Properties of Neural Networks.
In: \textit{International Conference on Learning Representations (ICLR)}, 2014.

\bibitem{Ribeiro2016}
Ribeiro, M.\,T.; Singh, S.; Guestrin, C.:
``Why Should I Trust You?'': Explaining the Predictions of Any Classifier.
In: \textit{Proceedings of the 22nd ACM SIGKDD International Conference on Knowledge Discovery and Data Mining},
S.~1135--1144, 2016.

\bibitem{Lipton2018}
Lipton, Z.\,C.:
The Mythos of Model Interpretability.
\textit{Queue}, 16(3), S.~31--57, 2018.

\bibitem{Marcus2018}
Marcus, G.:
Deep Learning: A Critical Appraisal.
\textit{arXiv preprint arXiv:1801.00631}, 2018.

\bibitem{Seshia2018}
Seshia, S.\,A.; Sadigh, D.; Sastry, S.\,S.:
Formal Specification for Deep Neural Networks.
\textit{arXiv preprint arXiv:1802.03023}, 2018.

\bibitem{Avizienis2004}
Avizienis, A.; Laprie, J.-C.; Randell, B.; Landwehr, C.:
Basic Concepts and Taxonomy of Dependable and Secure Computing.
\textit{IEEE Transactions on Dependable and Secure Computing}, 1(1), S.~11--33, 2004.

\bibitem{Leveson2011}
Leveson, N.\,G.:
\textit{Engineering a Safer World: Systems Thinking Applied to Safety}.
Cambridge, MA: MIT Press, 2011.

\bibitem{Varshney2017}
Varshney, K.\,R.; Alemzadeh, H.:
On the Safety of Machine Learning: Cyber-Physical Systems, Decision Sciences, and Data Products.
\textit{Big Data}, 5(3), S.~246--255, 2017.

\bibitem{Burton2017}
Burton, S.; Gauerhof, L.; Heinzemann, C.:
Making the Case for Safety of Machine Learning in Highly Automated Driving.
In: \textit{Computer Safety, Reliability, and Security (SAFECOMP)}.
Springer, 2017.

\bibitem{Huang2017}
Huang, X.; Kwiatkowska, M.; Wang, S.; Wu, M.:
Safety Verification of Deep Neural Networks.
In: \textit{Computer Aided Verification (CAV)}, S.~3--29, 2017.

\bibitem{Rumelhart1986}
Rumelhart, D.\,E.; Hinton, G.\,E.; Williams, R.\,J.:
Learning Representations by Back-propagating Errors.
\textit{Nature}, 323, S.~533--536, 1986.

\bibitem{LeCun2015}
LeCun, Y.; Bengio, Y.; Hinton, G.:
Deep Learning.
\textit{Nature}, 521, S.~436--444, 2015.

\bibitem{Goodfellow2016}
Goodfellow, I.; Bengio, Y.; Courville, A.:
\textit{Deep Learning}.
Cambridge, MA: MIT Press, 2016.

\bibitem{Dreossi2019}
Dreossi, T.; Donze, A.; Seshia, S.\,A.:
Compositional Falsification of Cyber-Physical Systems with Machine Learning Components.
\textit{Journal of Automated Reasoning}, 63, S.~1031--1053, 2019.

\bibitem{EUAI2019}
High-Level Expert Group on AI (Europäische Kommission):
\textit{Ethics Guidelines for Trustworthy AI}.
Brüssel: Europäische Kommission, 2019.

\bibitem{Zhao2020}
Zhao, X.; Nguyen, H.-D.; Bhatt, D.; Zulkernine, M.:
Assessing Safety of Machine Learning for Use in Safety-Critical Domains.
\textit{IEEE Transactions on Reliability}, 69(3), S.~1036--1056, 2020.

\bibitem{Kurd2003}
Kurd, Z.; Kelly, T.:
Establishing Safety Criteria for Artificial Neural Networks.
In: \textit{International Conference on Knowledge-Based Intelligent Information and Engineering Systems (KES)}.
Springer, 2003.

\bibitem{Koopman2019}
Koopman, P.; Fratrik, F.:
How Many Operational Design Domains, Objects, and Events?
In: \textit{Safety and Artificial Intelligence Workshop (AAAI)}, 2019.

\bibitem{Stocco2020}
Stocco, A.; Weiss, M.; Calzana, M.; Tonella, P.:
Misbehaviour Prediction for Autonomous Driving Systems.
In: \textit{42nd International Conference on Software Engineering (ICSE)}, 2020.

\bibitem{Pfeifer2021}
Pfeifer, H.; Herber, P.:
Challenges for Assuring AI-Based Safety Functions in Compliance with ISO~26262.
\textit{Proceedings of the 9th European Congress on Embedded Real Time Software and Systems (ERTS)}, 2021.

\bibitem{Doshi2017}
Doshi-Velez, F.; Kim, B.:
Towards a Rigorous Science of Interpretable Machine Learning.
\textit{arXiv preprint arXiv:1702.08608}, 2017.

\bibitem{Gauerhof2018}
Gauerhof, L.; Munk, P.; Burton, S.:
Structuring Validation Targets of a Machine Learning Function Applied to Automated Driving.
In: \textit{Computer Safety, Reliability, and Security (SAFECOMP)}.
Springer, 2018.

\bibitem{Sommerville2016}
Sommerville, I.:
\textit{Software Engineering}.
10.~Ausgabe. Harlow, UK: Pearson, 2016.

\bibitem{Wing2021}
Wing, J.\,M.:
Trustworthy AI.
\textit{Communications of the ACM}, 64(10), S.~64--71, 2021.

\bibitem{Tonella2022}
Tonella, P.:
When Testing AI Is (Not) Different.
\textit{arXiv preprint arXiv:2202.10524}, 2022.

\bibitem{Myllyaho2021}
Myllyaho, L.; Raatikainen, M.; M\"annis\"o, T.; Lepp\"anen, V.; Honkola, J.:
Systematic Literature Review of Validation Methods for AI Systems.
\textit{Journal of Systems and Software}, 181, 2021.

\end{thebibliography}

\end{document}
